\submittedby{Satyajit}

We use the usual convention for logical connectives: ``and'' has precedence over ``or''. Thus,
the condition defining the second family is interpreted as
\[ a\in V\quad\text{or}\quad\bigl(b\in V\text{ and }X\setminus V\text{ is finite}\bigr). \]

\textbf{Answer.} The collection $\T$ is \textbf{not a topology} on $X$. We show that it is not
closed under finite intersections.

\textbf{Counterexample.} Consider
\[ U=\{a,b\},\qquad V=X\setminus\{a\}. \]
We first verify that both sets belong to $\T$.

Since $a\in U$, the set $U$ belongs to the second family, and hence $U\in\T$.

Next, $b\in V$ and
\[ X\setminus V=\{a\}, \]
which is finite. Therefore $V\in\T$ as well.

However,
\[ U\cap V=\{a,b\}\cap(X\setminus\{a\})=\{b\}. \]
We now show that $\{b\}\notin\T$.

First, $\{b\}\notin\Pow(X\setminus\{a,b\})$, because $b\notin X\setminus\{a,b\}$. Also,
$a\notin\{b\}$, while $X\setminus\{b\}$ is infinite, since $X$ is infinite. Thus $\{b\}$ does
not satisfy
\[ a\in W\quad\text{or}\quad\bigl(b\in W\text{ and }X\setminus W\text{ is finite}\bigr) \]
either. Consequently, $\{b\}\notin\T$.

We have therefore found $U,V\in\T$ such that
\[ U\cap V=\{b\}\notin\T. \]
Hence $\T$ is not closed under finite intersections, and so
\[ \boxed{\T\text{ is not a topology on }X.} \]
\hfill$\square$

% REVIEW: only the "a or (b and cofinite)" reading is answered. Under the other grouping "(a in V or b in V) and X\V finite", T IS a topology (e.g. two cofinite sets meeting in neither a nor b give a subset of X\{a,b}); the remark below names that reading but does not settle it.
\begin{remark}[On the logical grouping]
If the intended condition had instead been
\[ (a\in V\text{ or }b\in V)\text{ and }X\setminus V\text{ is finite}, \]
then it would define a different collection. The counterexample above addresses the expression
as printed, using the standard precedence of ``and'' over ``or''.
\end{remark}
